\cabstractcn{
    此处格式已按模板设定，作者只需选择段落区域，输入替换之。模板中所有说明性文字用于注释格式与内容的要求，撰写论文时请删除。\\
    \hspace*{\parindent}中文摘要一般为300$\sim$400字，简要介绍毕业设计（论文）的研究目的、方法、结果和结论，语言力求精炼。中英文摘要均要有关键词，一般为3$\sim$8个。中英文摘要及关键词要相互对应。\\
    \hspace*{\parindent}中文摘要。“摘要”两字之间空一个全角空格或两个半角空格，字体为宋体二号字加粗，居中编排，摘要内容采用正文样式。中文关键词与摘要内容间隔一行，无缩进左对齐书写。“关键词：”采用宋体四号字加粗，关键词内容采用正文样式，且换行不缩进。关键词之间用逗号分隔。\\
    \hspace*{\parindent}英文摘要。此部分皆为Times New Roman字体。“ABSTRACT”为二号字加粗，居中编排。英文摘要内容采用正文样式，采用两端对齐。英文关键词与英文摘要内容间隔一行，无缩进左对齐书写。“KEY WORDS:”为四号字加粗，英文关键词采用正文样式，且换行不缩进，关键词之间用分号分隔，词义和中文关键词相同。“ABSTRACT”和“KEY WORDS”一律用大写字母，每个关键词的首字母要大写。
} %摘要这里由于使用了变量存储，所以段落间不要空格并且需要在最后添加\\强制换行并添加\hspace*{\parindent}缩进，这样很不优雅但是我也不像使用类似parbox宏包破坏代码的简洁性，不过摘要一般也不需要分段吧？

\ckeywordcn{
    关键词1，关键词2，关键词3，关键词4，关键词5，关键词6，关键词7
}

\cabstracten{
    [英文摘要应与中文摘要相对应。字体为Times New Roman小四号。]
}

\ckeyworden{
    Heat transfer；Mechanics；Combustion；Engine；Analysis
}

\makeabstract
